\section{Experimental Setup}

We design experiments to answer four research questions that map directly to our central claims about curation-contamination relationships:

\textbf{RQ1:} Does CCR vary systematically across filtering strategies, or are contamination rates independent of curation decisions?

\textbf{RQ2:} Are high-CCR examples \textit{causally necessary} for benchmark performance, or merely correlated with it?

\textbf{RQ3:} Does the Amplification Index (AI) distinguish filtering strategies' contamination effects on potentially contaminated versus clean benchmarks?

\textbf{RQ4:} Do contaminated high-influence examples exhibit structurally different influence patterns (IFR) than non-contaminated high-influence examples?

\subsection{Datasets}

\subsubsection{Training Corpus}

We use \textbf{RedPajama-Data-V2} (English, snapshot 2023-14) as our source corpus. RedPajama-V2 includes pre-computed quality signals including \texttt{ccnet\_perplexity} (Wikipedia LM perplexity scores), enabling reproducible filtering experiments without additional perplexity computation.

\begin{table}[h]
\centering
\begin{tabular}{@{}ll@{}}
\toprule
Property & Value \\
\midrule
Source & togethercomputer/RedPajama-Data-V2 \\
Language & English \\
Snapshot & 2023-14 \\
Quality signals & ccnet\_perplexity, ccnet\_bucket \\
Tokens per strategy & $\sim$1B (matched) \\
\bottomrule
\end{tabular}
\caption{Training corpus properties.}
\end{table}

\textbf{Why RedPajama-V2:} Pre-computed perplexity signals match production curation pipelines. Single snapshot eliminates temporal confounds.

\subsubsection{Evaluation Benchmarks}

\begin{table}[h]
\centering
\begin{tabular}{@{}lll@{}}
\toprule
Benchmark & Purpose & Examples \\
\midrule
MMLU & Primary contamination target & 14,042 \\
MMLU-Redux (2024+) & Time-stratified clean control & $\sim$500 \\
\bottomrule
\end{tabular}
\caption{Evaluation benchmarks.}
\end{table}

\textbf{Why MMLU:} Widely used benchmark with documented contamination concerns in web-scraped corpora.

\textbf{Why MMLU-Redux:} Post-training-cutoff benchmark for Amplification Index computation; assumed $<$0.01\% overlap with pre-2024 training data.

\subsection{Filtering Strategies}

We apply three filtering strategies to the same RedPajama-V2 source:

\begin{enumerate}
    \item \textbf{Perplexity-filtered}: Bottom 30\% by \texttt{ccnet\_perplexity} (low perplexity = high quality per Wikipedia LM)
    \item \textbf{Random-sampled}: Uniform random selection (baseline)
    \item \textbf{Inverse-perplexity}: Top 30\% by perplexity (control---expected low CCR)
\end{enumerate}

Each strategy produces a $\sim$1B token training corpus. The matched token budget isolates filtering effects from corpus size confounds.

\subsection{Model and Training}

\begin{table}[h]
\centering
\begin{tabular}{@{}ll@{}}
\toprule
Parameter & Value \\
\midrule
Model & Pythia-1B (GPT-NeoX architecture) \\
Parameters & 1.0B \\
Optimizer & AdamW \\
Learning rate & $2.5 \times 10^{-4}$ (cosine decay) \\
Warmup & 1\% of steps \\
Batch size & 512 $\times$ 2048 tokens \\
Weight decay & 0.1 \\
$\beta_1, \beta_2$ & 0.9, 0.95 \\
Gradient clipping & 1.0 \\
Seeds per strategy & 5 \\
\bottomrule
\end{tabular}
\caption{Training hyperparameters.}
\end{table}

Total training runs: 3 strategies $\times$ 5 seeds = 15 models.

\subsection{Evaluation Metrics}

\begin{table}[h]
\centering
\begin{tabular}{@{}lll@{}}
\toprule
Metric & Definition & Success Criterion \\
\midrule
CCR & Contamination contribution ratio & CCR(ppl) - CCR(rand) $>$ 0.1 \\
Degradation Ratio & $\frac{\Delta\text{acc(high-CCR)}}{\Delta\text{acc(random)}}$ & $\geq 1.5$ \\
Amplification Index & $\Delta$acc(MMLU) - $\Delta$acc(MMLU-Redux) & $> 0$, CI excludes zero \\
IFR & Influence fragility ratio & IFR(contam) $>$ IFR(non-contam) \\
\bottomrule
\end{tabular}
\caption{Evaluation metrics and success criteria.}
\end{table}

Statistical significance evaluated using bootstrap with 1000 resamples; we report 95\% confidence intervals.
